\chapter{The Quant Strats Role}
\label{ch:strats_role}
\chaptermark{The Quant Strats role}

What does a Strat actually do?  This chapter maps the sell-side
floor and walks through a typical day.  Part~VII grounds earlier
theory in the specific institution where you will work.  A Strat at
Goldman Sachs sits at the intersection of quantitative research,
trading, and engineering: you derive a pricing model from first
principles \emph{and} ship production code that runs it in real time
across every desk.  Understanding that --- hour by hour --- is
essential preparation for the off-cycle interview and your first
weeks on the floor.

\begin{backgroundbox}[Prerequisites]
From earlier in the book:
\begin{itemize}
  \item \Cref{ch:market_making}: the Avellaneda--Stoikov
    market-making model, reservation price, inventory skew, and the
    quoting loop.  A Strat calibrates and monitors models of exactly
    this kind.
  \item \Cref{ch:black_scholes}: Black--Scholes pricing, risk-neutral
    valuation, and the Greeks.  Morning recalibration of the
    volatility surface is the first task in a Strat's day.
  \item \Cref{ch:almgren_chriss}: optimal execution and the
    Almgren--Chriss framework.  Strats own the execution algorithms
    that implement these ideas in production.
\end{itemize}
\end{backgroundbox}

%% ================================================================
\section{The sell-side floor map}
\label{sec:strats_floor_map}
%% ================================================================

Walk onto the floor of a bulge-bracket bank (Goldman Sachs, Morgan
Stanley, J.P.~Morgan) and you see a single open-plan room the size
of a football pitch: rows of desks, six or more monitors each, loud
under fluorescent light.  The room is organised not by seniority but
by \emph{product} --- clusters for equities, rates, credit, FX,
commodities.  Within each cluster, five roles sit close enough to
shout at each other.  Proximity is the primary communication
channel: when a trader needs a price in ten seconds, email will not
do; a shout across three desks will.

\subsection{Product desks}

A \textbf{product desk} is a self-contained business unit organised
around a related set of tradeable instruments.  The major desk
families at Goldman Sachs:

\begin{itemize}
  \item \textbf{Rates} --- government bonds, interest-rate swaps,
    swaptions, caps, floors, inflation-linked products.
  \item \textbf{Equities} --- single stocks, equity indices, equity
    derivatives (calls, puts, exotics, variance swaps), ETFs, program
    trading.
  \item \textbf{Credit} --- corporate bonds, credit default swaps
    (CDS), CDS indices, collateralised debt obligations (CDOs),
    leveraged loans.
  \item \textbf{Foreign exchange (FX)} --- spot, forwards, FX swaps,
    FX options (vanillas and barriers), cross-currency basis swaps.
  \item \textbf{Commodities} --- oil, natural gas, metals,
    agriculture, power, emissions.  Each sub-commodity often has its
    own desk.
\end{itemize}

Each desk has its own P\&L, risk limits, and regulatory capital
allocation.  A junior hire is assigned to one desk for two to three
years; the desk determines your first asset class, your models, and
your traders.

\subsection{The five roles}

\paragraph{Traders.}
The trader owns the risk book: which positions to take, which
client axes to fill, when to hedge.  (An \emph{axis} is a standing
interest --- ``the desk is a buyer of 5Y protection'' --- that Sales
circulates to clients.)  Traders are measured on P\&L ---
the single number determining bonus, reputation, and survival.  A
desk head manages a book with gross notional in the billions and a
daily P\&L swinging by several million.  Traders speak in terse
shorthand, need answers \emph{now}, and are the primary internal
client for every other role.

Two archetypes.  \textbf{Flow traders} (``market makers'' on the
sell side) quote prices, manage inventory, and earn the bid-ask
spread.  \textbf{Proprietary traders} --- rare since the Volcker Rule
(2010) --- take directional bets with the firm's capital.  Most
modern bank trading is flow, so the Strat's job is to help price,
hedge, and risk-manage \emph{client-driven} positions.

\paragraph{Sales.}
Sales sits between the trader and the external client (hedge funds,
pension funds, corporates, sovereigns), translating a client's
hedging or investment need into a trade request for the trader to
price.  Sales owns the client relationship, not risk; compensation
tracks client revenue, not P\&L.  A good salesperson explains Greeks
to a client treasurer but does not derive them.

\paragraph{Quantitative researchers (QR).}
QR publishes models: new stochastic processes for underlyings,
better calibration algorithms, novel risk metrics.  They write
papers and prototypes (Python or MATLAB) and hand models to Strats
for productionisation.  A QR researcher might spend six months
deriving a local-stochastic-vol calibration; the Strat turns it
into code running in milliseconds across every desk by next
quarter.  At Goldman Sachs, QR is a separate division from
Engineering and Strats, though all three share the floor.

\paragraph{Strats (strategists).}
The Strat (``strategist'') is the bridge: writes production code
\emph{and} has the depth to derive the models the code implements.
This dual requirement distinguishes the role from software
engineering (infrastructure without model logic) and from QR (models
without production code).  Goldman Sachs formalised the title in
the 1990s; other banks call the same role ``desk quant,'' ``quant
developer,'' or ``front-office quant.''

\paragraph{Engineering (Technology).}
Engineering builds the platforms everything runs on: networking,
databases, messaging, trade capture, regulatory reporting, the
firm-wide risk aggregation engine.  Engineers own reliability,
throughput, and internal developer experience, not model logic.
The off-cycle placement sits in Engineering --- ``at the critical
center of our business,'' since every trade, risk number, and
client report flows through infrastructure Engineering maintains.

\begin{keyfactbox}[Floor map summary]
\begin{tabular}{llll}
\toprule
\textbf{Role} & \textbf{Owns} & \textbf{Measured by} & \textbf{Primary output} \\
\midrule
Trader   & Risk book       & P\&L            & Positions, hedges \\
Sales    & Client relationship & Client revenue & Trade requests \\
QR       & Model research   & Publication, adoption & Papers, prototypes \\
Strat    & Production models & Model accuracy, uptime & Production code \\
Eng.     & Infrastructure   & Reliability, throughput & Platforms, tooling \\
\bottomrule
\end{tabular}
\end{keyfactbox}

\subsection{How the roles interact}

A client trade illustrates the roles.  A pension fund calls Sales:
``We want to hedge our equity exposure for six months --- price us
a collar on the S\&P~500 with a 95\% put and a 105\% call.''  Sales
passes it to the Trader, who asks the Strat: ``Run the collar
through the engine --- fair value and vega?''  The Strat evaluates
in SecDB, which calls the vol-surface calibration code the Strat
wrote last quarter, computes price and Greeks, and returns the
result.  The Trader adds a margin for execution risk and target
return; Sales quotes the client.  If the client trades,
Engineering's trade-capture system books it, SecDB updates risk in
real time, and the overnight batch (Strat code on Engineering
infrastructure) produces the regulatory report.

\begin{pitfallbox}[Strat $\neq$ Software Engineer]
``Strat'' and ``Software Engineer'' are not interchangeable.  A bank
software engineer builds generic infrastructure (messaging, databases,
UIs) and is not expected to derive a pricing formula.  A Strat
implements a \emph{specific} pricing or risk model, owns its
mathematical correctness, and sits physically on the trading desk
rather than in a separate technology building.  The Engineering job
description emphasises ``risk management, big data, mobile'' --- all
infrastructure.  The Strat job description (a separate hiring track)
emphasises ``pricing models, Greeks, calibration'' --- the model
layer.  Know which track you are interviewing for.
\end{pitfallbox}

%% ================================================================
\section{Strats as the bridge}
\label{sec:strats_bridge}
%% ================================================================

The Strat role sits at the intersection of three competency axes.  A
Strat must be able to:

\begin{enumerate}
  \item \textbf{Derive}: read a QR paper, reproduce the derivation,
    identify which approximations matter, decide whether the model is
    fit for purpose on a given desk.
  \item \textbf{Implement}: write production code --- the firm's
    proprietary language (Slang at Goldman Sachs;
    \cref{sec:strats_secdb}) or C++/Python --- that runs in real
    time, handles edge cases, and degrades gracefully under stress.
  \item \textbf{Communicate}: explain to the trader, in thirty
    seconds, why the model gives a particular price, what changed
    overnight, and what the risk is if the model is wrong.
\end{enumerate}

No other role requires all three.  Traders communicate and decide
but do not derive or implement; QR derives but does not own
production code; Engineers implement but do not derive.  Only the
Strat can trace a suspicious P\&L from the trader's screen back
through production code to the mathematical model and identify
whether the bug is in the formula, the implementation, or the input.

\subsection{The translation problem}

A Strat's value is most visible when something goes wrong.  A
scenario that recurs monthly on any active desk: the trader looks
at the morning risk report: ``My vega is \$500K per vol point on
this swaption position, but yesterday it was \$300K.  Nothing
changed overnight.  What happened?''

An engineer cannot answer: it requires understanding the model ---
did the surface recalibrate to a different shape, changing vega?  A
QR researcher might answer in principle but cannot trace the bug
through production code to the line where calibration parameters
are read.  The Strat can: check whether the surface moved (yes --- a
new quote shifted SABR $\rho$), verify the calibration code ingested
the quote correctly (yes), and explain the vega change is real, not
a bug, and the position is now more sensitive to vol moves than
yesterday.

End-to-end debugging --- market data through model mathematics
through code to the screen --- is the Strat's defining skill, and
the hardest to develop, requiring simultaneous fluency in three
domains taught separately at university: mathematics, computer
science, and finance.

\subsection{Daily responsibilities}

A Strat's daily work splits into five recurring categories, each
drawing on earlier theory.

\paragraph{1.\ Model calibration.}
Every morning the vol surface (\cref{ch:vol_surface}) is
recalibrated to the previous close and overnight moves.  The Strat
owns the calibration code: SABR to the swaption grid (rates),
local-vol or Heston to the equity options chain, hazard-rate curve
to CDS spreads (credit).  When the fit breaks --- a new strike
outside the interpolation range, an optimiser that fails to
converge --- the Strat is paged first.

Calibration is not one-off.  The surface drifts as new trades print.
A good routine runs incrementally (updating only parameters affected
by new data), handles missing data (some strikes have no fresh
quotes), and produces an \emph{arbitrage-free} surface: no butterfly
with negative price, no calendar spread violating time-monotonicity.
Enforcing these under time pressure is a real engineering challenge
belonging to the Strat.

\paragraph{2.\ Risk computation.}
Greeks (\cref{ch:greeks_hedging}) are recomputed for the full book
start-of-day and updated intraday.  The Strat owns the risk
engine's model layer: which bumps to apply (absolute or relative,
and how large), how to handle cross-Greeks (e.g.\ \emph{vanna},
delta's sensitivity to vol), and which approximations are acceptable
(closed-form where possible, Monte Carlo where necessary,
finite-difference as fallback).  The output feeds the trader's risk
dashboard and the firm-wide report to management and regulators.

A subtlety: on-screen Greeks must be \emph{consistent} with the
price.  If price comes from Monte Carlo but delta from a closed-form
approximation, the trader's hedge is wrong.  Ensuring consistency
(all Greeks from the same pricing function, ideally via AD;
\cref{sec:strats_secdb}) is a core Strat responsibility.

\paragraph{3.\ Trade pricing.}
When Sales brings a structured-product request --- say, an
autocallable on the EuroStoxx~50 --- the trader asks the Strat for
fair value.  The Strat runs the payoff through the pricing engine,
checks Greeks, returns a price.  If no existing model handles the
payoff, the Strat builds one under time pressure (the client wants
a price by end of day).

Structured products are where mathematical depth is directly
monetised.  An autocallable has path-dependent payoffs depending on
the joint dynamics of underlying, rates, and dividends.  Correct
pricing needs the right model (local vol, stochastic vol,
local-stochastic vol), calibration to the current market, Monte
Carlo with enough paths for precision, and Greeks by pathwise
differentiation or likelihood-ratio methods.  The Strat does all of
this; the trader relies on the number.

\paragraph{4.\ Prototyping and production.}
New models start as prototypes: a Jupyter notebook testing a QR idea
against historical data.  If it works, the Strat rewrites it in
production code, integrates with the firm's pricing/risk database
(SecDB at Goldman Sachs; \cref{sec:strats_secdb}), and deploys.  The
Strat then monitors, fixes bugs, and tunes parameters as conditions
change.

Prototype-to-production is where many models die.  A notebook model
that works on clean data may fail catastrophically in production on
missing data, stale quotes, corporate actions (splits, dividend
changes), or extreme moves.  The Strat anticipates these failures,
builds guards, and ensures graceful degradation (fall back to
yesterday's calibration) rather than catastrophic (price of zero,
automatic hedge triggered).

\paragraph{5.\ P\&L explanation.}
Each morning the trader asks: ``Why is my P\&L what it is?''  The
Strat decomposes overnight P\&L into contributions from delta, gamma,
vega, theta, and unexplained residual.  This decomposition
(\cref{ch:greeks_hedging}) is not just accounting --- a large
residual signals model mis-specification, and the Strat must
diagnose it.

The standard decomposition is a Taylor expansion of the portfolio
value $V$:
\begin{equation}
  \Delta\text{P\&L}
  \;\approx\;
  \underbrace{\Delta \cdot \Delta S}_{\text{delta}}
  + \underbrace{\tfrac{1}{2}\Gamma\,(\Delta S)^2}_{\text{gamma}}
  + \underbrace{\mathcal{V}\cdot\Delta\sigma}_{\text{vega}}
  + \underbrace{\Theta\cdot\Delta t}_{\text{theta}}
  + \underbrace{\text{residual}}_{\text{unexplained}},
  \label{eq:strats_pnl_explain}
\end{equation}
where $\Delta S$ is the change in underlying, $\Delta\sigma$ the
change in implied volatility, and $\Delta t$ the elapsed time.  A
``clean'' P\&L explain has residual below 5\% of total; above 10\%
triggers an investigation.

\begin{intuitionbox}
The Strat translates between two languages: mathematics (QR) and
risk/money (trader).  Fluent in both, the Strat is the only person
on the floor who can debug a pricing discrepancy end to end --- from
the SDE to the C++ inner loop to the number on the trader's screen.
\end{intuitionbox}

\begin{keyfactbox}[The five daily tasks]
\begin{enumerate}
  \item \textbf{Calibrate}: fit the vol surface / curve to today's market.
  \item \textbf{Compute risk}: rerun Greeks for the full book.
  \item \textbf{Price}: value new trades and structured products.
  \item \textbf{Build}: prototype new models, ship to production.
  \item \textbf{Explain}: decompose P\&L into Greek contributions.
\end{enumerate}
Every other task is a variant or combination of these five.
\end{keyfactbox}

%% ================================================================
\section{A day in the life}
\label{sec:strats_day_in_life}
%% ================================================================

A composite sketch of a typical day for a junior Strat on an equity
derivatives desk.  Specifics vary by desk, seniority, and firm, but
the rhythm is representative.

\subsection{7:00 AM --- Arrive, check overnight}

Asia closed hours ago; Europe is about to open.  The Strat checks
the overnight batch: did end-of-day risk complete?  Any calibration
failures?  The dashboard (built by Engineering, populated by Strat
code) is green everywhere except one: a single-stock option's vol
surface failed to converge because the stock gapped 5\% overnight
on earnings.  The Strat opens the calibration log, identifies the
problematic strike, widens optimiser bounds, reruns, verifies
sensible Greeks, and pushes the fix before the trader arrives at
7:15.

\subsection{7:30 AM --- Morning meeting}

The desk head runs a five-minute stand-up.  The trader: ``Long
gamma on the index, short gamma on single names.  Correlation book
made money overnight --- realised correlation dropped.  Asia was
quiet.''  The Strat reports the calibration fix, flags that gamma on
the at-the-money straddle is unusually high (the earnings gap
compressed the surface near the money), and recommends rehedging.
Sales mentions a client wanting a six-month barrier option price by
noon.

The morning meeting tests the Strat's communication.  The trader
does not want to hear about Levenberg--Marquardt failing to
converge; the trader wants: ``The surface on that name is steeper
near the money, raising gamma exposure by 20\%.  To keep the same
risk profile, sell about 500 ATM-straddle contracts.''  Translating
model output into actionable trading language in real time takes
months to develop.

\subsection{8:00 AM --- P\&L explain}

The Strat runs the P\&L decomposition using
\cref{eq:strats_pnl_explain}.  Yesterday's P\&L was \$+1.2M.  The
decomposition:

\begin{center}
\begin{tabular}{lr}
\toprule
\textbf{Component} & \textbf{Contribution} \\
\midrule
Delta (book was net long, market rallied) & $+\$800\text{K}$ \\
Theta (time decay on short options)       & $+\$300\text{K}$ \\
Vega (implied vol ticked up slightly)     & $+\$50\text{K}$ \\
Cross-gamma and higher-order              & $+\$20\text{K}$ \\
Unexplained residual                      & $+\$30\text{K}$ \\
\midrule
\textbf{Total}                            & $+\$1{,}200\text{K}$ \\
\bottomrule
\end{tabular}
\end{center}

The residual is 2.5\% --- well within the 5\% threshold.  The Strat
attributes it to bid-ask crossing on illiquid positions and rounding
in finite-difference Greek bumps.  The trader signs off.  Had the
residual been \$150K (12.5\%), the Strat would need to investigate:
mis-specified model?  Stale market-data feed?  Mis-booked position?

\subsection{10:00 AM --- Price the barrier option}

Sales has formalised the request: a six-month down-and-out put on
an equity index, barrier at 85\% of spot, \$50M notional --- a
path-dependent exotic that knocks out worthless if the index
touches the barrier in six months.  The Strat prices via the
local-vol engine (100,000 Monte Carlo paths, Euler discretisation,
500 time steps), sanity-checks against Heston (agreement within
10~bps of notional), computes Greeks by pathwise differentiation,
and sends a one-pager: fair value, delta, gamma, vega, theta,
knock-out probability (18\%), and ``barrier delta'' (sensitivity to
the barrier level).  The trader adds a 50~bp margin and Sales
quotes.

\subsection{12:00 PM --- Code review}

A fellow Strat has submitted a PR refactoring dividend-handling
in the equity pricing engine.  The change affects how discrete
dividends are interpolated between ex-dates, shifting the forward
curve and every option price on the desk.  The reviewer checks
mathematical correctness (does interpolation preserve no-arbitrage?),
numerical stability (what if two ex-dates are one day apart?), and
test coverage (do regressions catch a pricing change above 1~bp
on the benchmark book?).  Model-code review is more demanding than
infrastructure review: a subtle math error mis-hedges the desk and
costs real money.

\subsection{2:00 PM --- Fix a production bug}

A neighbouring-desk trader reports negative vega on a long
swaption --- clearly wrong.  The bug: a sign error in the SABR vega
bump, introduced two weeks ago when a QR patch switched from
absolute to relative bumps.  The absolute bump $\sigma \to \sigma +
\epsilon$ became $\sigma \to \sigma(1 + \epsilon)$, but read-back
still subtracts $\epsilon$ rather than dividing by $(1 + \epsilon)$.
The Strat writes a fix, adds a regression test for sign of vega on
a vanilla long call (always positive), gets a second-Strat review,
and pushes.  The corrected vega appears within minutes.

\begin{pitfallbox}[Model bugs cost real money]
A Greek sign error is not an academic mistake.  A trader seeing
negative vega on a long option might \emph{sell} vol to ``hedge,''
which \emph{doubles} their vega.  On a \$500K-per-vol-point
position, a 2-point vol move then costs \$1M instead of \$500K.
Every model bug has a dollar cost, and the Strat who ships it owns
the consequences.  Code review, regression testing, and careful
deployment are the primary defence against P\&L-destroying errors.
\end{pitfallbox}

\subsection{5:00 PM --- End of day}

Europe closes.  The Strat triggers the end-of-day batch: full recalibration, full
risk, regulatory report (about 45 minutes --- tens of thousands of
trades, each requiring a pricing evaluation).  While it runs, the
Strat reviews a QR paper on a stochastic local-vol model that
might improve barrier-option pricing; the key result (a PDE for
the local-vol component conditional on the stochastic-vol state)
is tractable and implementable with a moderate refactor.
Prototype notes go in the backlog.

\subsection{6:00 PM --- Batch monitoring and wrap-up}

The batch completes.  The Strat checks output (all calibrations
converged, all risk numbers within bounds, no limit breaches),
queues the regulatory report, commits the day's changes, updates
the tracker, and leaves --- 6:00--6:30~PM on most days.  On bad days
(batch-time model failure, regulatory deadline, product launch) it
runs much later.

\begin{pitfallbox}[The Strat is not a 9-to-5 role]
A Strat's day is bounded by the market, not the clock.  A model
breaking in Asian hours with live desk risk means a 3:00~AM page;
a new product launch means weekend work on the pricing model.
Quarter- and year-end add pressure: every position revalued, every
model validated, every regulatory report filed.  The role demands
deep technical skill \emph{and} the operational resilience to be
the person called when something goes wrong at an inconvenient hour.
\end{pitfallbox}

%% ================================================================
\section{Career arcs}
\label{sec:strats_career}
%% ================================================================

The standard Goldman Sachs progression (and most bulge-brackets')
follows a well-defined hierarchy.  Understanding it matters for
interview preparation and career planning.

\paragraph{Analyst / Off-cycle (0--3 years).}
You join as analyst (or off-cycle intern who converts).  Assigned
to one desk (rates, equities, credit, FX, commodities), you spend
year one learning the desk's products, models, and codebase.
Primary output: bug fixes, small features, the daily
calibration/risk cycle (\cref{sec:strats_day_in_life}).  The
Engineering off-cycle posting specifies ``real responsibilities
alongside fellow interns and our people'' --- you write production
code from week one, not observe.

Year one is steep: Slang, SecDB, the desk's product mathematics,
and the floor's unwritten conventions.  The fastest way to earn
trust is to fix a bug that has annoyed the trader for months ---
proof you can read the codebase, understand the model, and ship a
correct fix under pressure.

\paragraph{Associate (3--5 years).}
You own a specific model or product --- the person the trader calls
when it misbehaves.  You lead small projects: replacing a
calibration, building a pricing model for a new product, migrating
legacy code.  The firm expects you to identify model deficiencies
proactively, not just react to trader complaints.

\paragraph{Vice President (5--10 years).}
You own a broader area: the vol-surface calibration stack, or the
risk engine for a product suite.  You manage junior Strats, set
technical direction, and interact with senior traders and desk
heads on model choices.  Judgment calls: is this model good enough
or does it need replacing?  Is the engineering cost of a refactor
justified by the pricing improvement?  VPs also represent the desk
in firm-wide model-governance reviews, where senior risk managers
challenge assumptions and demand evidence of robustness.

\paragraph{Managing Director (10+ years).}
An MD Strat is rare, running a full Strats group (``Global Rates
Strats,'' ``Equity Derivatives Strats'').  The role is as much
people management, hiring, and strategy as code and models: an MD
decides which models to invest in, which platforms to build, and
how to allocate 10--30 Strats across competing priorities.

\begin{keyfactbox}[Career ladder summary]
\begin{tabular}{lll}
\toprule
\textbf{Level} & \textbf{Typical tenure} & \textbf{Scope} \\
\midrule
Analyst   & 0--3 years  & Individual tasks, bug fixes, daily cycle \\
Associate & 3--5 years  & Owns one model or product \\
VP        & 5--10 years & Owns a product suite, manages juniors \\
MD        & 10+ years   & Runs a Strats group, sets strategy \\
\bottomrule
\end{tabular}
\end{keyfactbox}

\subsection{Exit paths}

The sell side is a powerful launchpad.  After three to five years
a Strat's skillset --- production coding, model derivation, market
intuition, operational discipline --- is highly valued outside the
bank.

\begin{itemize}
  \item \textbf{Hedge funds.} The most common exit.  Multi-strategy
    funds (Citadel, Millennium, Point72) hire experienced Strats to
    run systematic desks or build pricing infrastructure.  Direct
    skill transfer --- pricing, Greeks, calibration, production
    coding --- and pay at a top fund can be 2--5$\times$ the bank
    at the same seniority.
  \item \textbf{Proprietary trading firms.} Jane Street, Two Sigma,
    Jump hire Strats for quantitative depth and production
    discipline; low-latency experience is especially valued.
  \item \textbf{Technology.} Some Strats move to big tech (Google,
    Meta, Amazon) where systems-thinking and low-latency skills
    transfer.  Pay is competitive (below hedge-fund levels) and
    work-life balance is better.
  \item \textbf{Quantitative research.} A smaller number move to
    pure QR --- same bank or a fund --- leaving production to
    others.  Suits Strats who prefer derivation to debugging.
  \item \textbf{Fintech and startups.}  A growing path: Strats bring
    credibility and domain expertise to trading platforms, risk
    systems, and crypto infrastructure.
\end{itemize}

\subsection{What Goldman Sachs looks for in off-cycle hires}

The Engineering posting emphasises three themes that apply equally
to Strats candidates:

\begin{enumerate}
  \item \textbf{Problem-solving under constraints.}  The posting
    calls for ``innovative strategic thinking and immediate, real
    solutions.''  Translation: given a problem, can you decompose
    it, find the binding constraint, and produce a working solution
    quickly?  Probability puzzles and live coding probe this.
  \item \textbf{Adaptability.}  ``Evolve, adapt to change and thrive
    in a fast-paced global environment.''  The sell side churns:
    new regulations, products, models.  The interviewer wants
    evidence you have operated in ambiguous, shifting environments.
  \item \textbf{Collaboration.}  ``Creative collaborators'' is the
    exact phrase.  The Strat role is collaborative by construction
    --- you work with traders, QR, sales, and engineering daily.
    The interviewer looks for signals that you can communicate
    technical ideas to non-technical stakeholders and vice versa.
\end{enumerate}

\begin{intuitionbox}
The single most important thing Goldman Sachs wants in a junior hire
is \emph{evidence of building things that work}: a side project, an
open-source contribution, a competition result, a research
implementation --- anything that shows you can go from idea to
working code under constraints.  IMC Prosperity is exactly this: it
shows you can design, implement, and iterate on a trading algorithm
under time pressure with real (simulated) consequences.
\end{intuitionbox}

%% ================================================================
\section{Technology: SecDB and Slang}
\label{sec:strats_secdb}
%% ================================================================

Goldman Sachs runs one of the most sophisticated proprietary
platforms in finance.  Two components are central to a Strat's daily
work: \textbf{SecDB} and \textbf{Slang}.

\subsection{SecDB: the pricing and risk database}

SecDB (\textbf{Sec}urities \textbf{D}ata\textbf{B}ase) is Goldman
Sachs's firm-wide pricing, risk-management, and position-management
system.  Built in the early 1990s and continuously developed since,
SecDB is not a relational database: it is a \emph{graph-based
compute engine} representing every trade, instrument, market-data
input, and pricing model as a node in a directed acyclic graph
(DAG).  When a market-data input changes --- a spot, a vol quote, a
yield-curve point --- SecDB propagates the change, recomputing
every dependent price and risk number automatically.

\begin{keyfactbox}[SecDB in one sentence]
SecDB is a reactive DAG-based compute engine that prices every trade
in the firm, computes every Greek, and propagates every market-data
change in near-real time.
\end{keyfactbox}

Scale is extraordinary.  At peak, SecDB holds pricing models for
hundreds of thousands of live trades across every asset class.
Every desk (rates, equities, credit, FX, commodities) uses the same
instance, so a single SABR-calibration change affects every desk
using SABR.  The benefit is consistency (one model, one price, one
risk number per trade); the cost is global bug propagation --- hence
the discipline of code review and testing.

The DAG structure shapes how a Strat thinks about pricing.  Leaves
are \textbf{market-data nodes}: spots, yield-curve points, vol-surface
quotes, dividend forecasts, updated in real time.  Above are
\textbf{model nodes}: calibration routines producing parameters
(SABR, Heston, hazard rates).  Above those, \textbf{trade nodes}:
one per trade, with edges to the model and market-data nodes it
depends on.  A market-data update triggers an upward traversal,
recomputing every dependent node.  Every price and Greek is always
consistent with the latest data, without manual intervention.

\begin{keyfactbox}[The SecDB graph]
\begin{center}
\begin{tabular}{ll}
\textbf{Layer} & \textbf{Contents} \\
\midrule
Trade nodes     & One per trade: price, Greeks, P\&L \\
Model nodes     & Calibration outputs: SABR params, Heston params, \ldots \\
Market-data nodes & Spot, yield curve, vol quotes, dividends \\
\end{tabular}
\end{center}
A change at any lower layer propagates upward automatically.
The Strat writes the model-node logic; Engineering writes the
market-data and infrastructure layers.
\end{keyfactbox}

\subsection{Slang: the pricing language}

\textbf{Slang} (\textbf{S}ecDB \textbf{Lang}uage) is Goldman Sachs's
proprietary functional language for writing pricing and risk models
inside SecDB: statically typed, garbage-collected, functional, with
built-in automatic differentiation (AD) --- the same technique
behind modern ML frameworks, applied here to compute Greeks.

\begin{intuitionbox}
AD in Slang means that when a Strat writes a pricing function, the
system computes every Greek (delta, gamma, vega, theta, rho, and
arbitrary cross-Greeks) \emph{automatically}, with no differentiation
code.  This eliminates an entire class of bugs: Greeks are exactly
consistent with the pricing function by construction.
\end{intuitionbox}

Slang's design features explain why Goldman Sachs built a custom
language instead of using C++ or Python:

\begin{enumerate}
  \item \textbf{Automatic differentiation (AD).}  Every Slang
    function differentiates automatically with respect to any input.
    Greeks are computed by the runtime, exact to machine precision,
    consistent with the price by construction, zero extra code.
  \item \textbf{Lazy evaluation.}  Slang evaluates expressions only
    when results are needed.  In SecDB's DAG a market-data change
    triggers recomputation only of nodes actually depending on the
    input, not the whole graph --- what makes near-real-time risk at
    scale feasible.
  \item \textbf{Serialisability.}  Every Slang object serialises to
    disk and back, letting SecDB checkpoint and restore the pricing
    graph --- critical for end-of-day batches, disaster recovery,
    and regulatory audits.
  \item \textbf{Type safety.}  The static type system catches common
    errors --- passing a volatility where a price was expected,
    confusing a date with a number --- at compile time.
\end{enumerate}

The Slang VM runs on every desk's compute infrastructure.  A price
request routes to a Slang process that evaluates the pricing graph
and returns the result.  The VM is optimised for two workloads:
single-trade pricing (latency-critical, for live quoting) and batch
risk (throughput-critical, for end-of-day and regulatory reports).

A junior Strat's first months are spent learning Slang and the
SecDB graph model.  The curve is steep --- Slang is not taught at
any university and has no public documentation --- but the payoff
is that a fluent Strat can implement, test, and deploy a new
pricing model in days rather than the weeks or months a conventional
C++ codebase would take.

\begin{pitfallbox}[Slang is not open-source]
You cannot practise Slang before joining Goldman Sachs: entirely
proprietary, no public compiler, no published textbook.  What you
\emph{can} practise is functional programming (Haskell, OCaml,
Python's functional idioms), AD (JAX, PyTorch autograd), and
graph-based computation (TensorFlow dataflow, Dask).  These
transferable skills accelerate the Slang transition, and in
interviews, comfort with functional programming and AD is a strong
positive signal.
\end{pitfallbox}

\subsection{Comparable systems at other banks}

Goldman Sachs is not alone.  The major systems, each serving the
same purpose (unified pricing, risk, position management) but
differing in language, architecture, and maturity:

\begin{center}
\begin{tabular}{llll}
\toprule
\textbf{Bank} & \textbf{System} & \textbf{Language} & \textbf{Architecture} \\
\midrule
Goldman Sachs  & SecDB   & Slang (proprietary) & Reactive DAG \\
J.P.~Morgan    & Athena  & Python / C++        & Grid compute \\
Morgan Stanley & In-house & C++ / C\#           & Service-oriented \\
Barclays       & BARX    & Java / C++          & Service-oriented \\
\bottomrule
\end{tabular}
\end{center}

The skills you build on \emph{any} of these systems --- reactive
computation, AD, model calibration, risk aggregation --- transfer
directly.  The specific language is secondary; the design patterns
are universal.

\subsection{Why SecDB matters for your career}

SecDB is widely regarded as one of Goldman Sachs's most significant
competitive advantages: unified risk across all desks, consistent
pricing of complex structured products, rapid deployment of new
models.  Its decades-long head start and Slang's expressiveness have
kept it at the frontier.  As a Strat you are both user and builder;
understanding the architecture is not optional.

The SecDB/Slang ecosystem teaches two things hard to learn
elsewhere.  First, writing code where \emph{correctness is the
primary constraint}: in most software a bug means a bad user
experience; in pricing code it means lost money.  The discipline
of provably correct financial code (comprehensive tests, formal
review, staged deployment) is highly valued by hedge funds and
prop shops.  Second, how a large-scale reactive system works:
propagating changes through a graph, handling stale data,
balancing latency against throughput --- problems in every
real-time data system, from ad platforms to autonomous vehicles.

\Cref{ch:secdb_slang} covers SecDB's architecture and Slang's
language features in depth.

%% ================================================================
%% Required boxes
%% ================================================================

\begin{prosperitybox}
Prosperity's \texttt{Trader} class --- \texttt{run()} called once
per iteration --- is a miniature Strat event loop.  Each call
receives the market state (\texttt{TradingState}) and returns
orders, conversions, and serialised state, just as a Strat's
pricing code must return price and Greeks within a single SecDB
evaluation.  The constraint that all state live in
\texttt{traderData} (a string persisting between iterations)
mirrors SecDB's requirement that model state be serialisable and
reproducible.  A clean, stateful \texttt{Trader.run()} handling
edge cases is practice in the core Strat discipline:
deterministic, stateful computation under time pressure with no
room for unhandled exceptions.
\end{prosperitybox}

\begin{gsbox}
\textbf{Interview prep for the Strat and Engineering tracks.}
Goldman Sachs quantitative interviews have four components:
\begin{enumerate}
  \item \textbf{Probability and brainteasers.}  Classics: ``Expected
    coin flips to get two heads in a row?''  ``100 balls, 50 red and
    50 blue, drawn without replacement --- probability the last is
    red?''  Tests reasoning under uncertainty and structured
    communication.
  \item \textbf{Black--Scholes and Greeks.}  Derive Black--Scholes
    from risk-neutral pricing (\cref{ch:black_scholes}), explain
    why ATM call delta is approximately 0.5, sketch gamma vs.\
    spot.  Strat interviews go deeper here than Engineering.
  \item \textbf{Coding test.}  Timed online assessment or live
    whiteboard: data structures and algorithms at strong-undergrad
    level, Python or C++.
  \item \textbf{Market intuition.}  ``If the Fed raises rates by
    25~bps, what happens to a 10-year Treasury?  To equities?
    Why?''  Tests connecting theory to real market moves.
\end{enumerate}
\Cref{ch:interview_problems} contains a bank of practice problems
covering all four categories.
\end{gsbox}

\begin{historybox}
Emanuel Derman's memoir \emph{My Life as a Quant} (2004) is the
canonical account of quantitative finance's early days on Wall
Street \citeDerman{}.  Derman, a theoretical physicist who joined
Goldman Sachs in 1985, helped build the firm's first-generation
options-pricing models and shaped the culture that became the
modern Strats group.  His description of the intellectual
excitement --- and operational chaos --- of building pricing models
in the late 1980s remains the best introduction to what the Strat
role \emph{feels like} from inside.  Two themes.  First, models are
always wrong; the question is whether they are useful enough, and
the Strat's judgment on ``useful enough'' is what the trader pays
for.  Second, the half-life of a pricing model is short --- markets
evolve, products change, and the Strat who stops learning stops
being useful.  Derman's arc --- physicist to junior quant to MD to
academic --- is a template in \cref{sec:strats_career}.
\end{historybox}

%% ================================================================
\section*{Key facts to memorise}
\label{sec:strats_key_facts}
%% ================================================================

\begin{itemize}
  \item The sell-side floor is organised by \textbf{product}, not
    function.  Traders, sales, QR, Strats, and engineering sit
    together per product area.
  \item A \textbf{Strat} writes production code \emph{and} derives
    the models it implements --- the dual competency separating the
    role from software engineering (code without models) and QR
    (models without production code).
  \item Daily loop: morning calibration and P\&L explain, mid-day
    trade pricing and prototyping, afternoon bug fixes and
    production pushes, evening reading.
  \item \textbf{SecDB} is Goldman Sachs's DAG compute engine for
    pricing and risk.  Every trade prices through it; market-data
    changes propagate through the graph automatically.
  \item \textbf{Slang} is the proprietary functional language for
    writing pricing models.  Built-in AD computes all Greeks for
    free, eliminating a major class of implementation bugs.
  \item Progression: Analyst $\to$ Associate $\to$ VP $\to$ MD.
    Common exits: hedge funds, prop shops, tech, pure QR.
  \item Off-cycle interviews test four axes: probability puzzles,
    Black--Scholes and Greeks, coding, market intuition.
  \item Prosperity's \texttt{Trader.run()} is a miniature Strat
    event loop: receive state, compute, return orders, serialise.
    Clean, stateful, deterministic code under time pressure --- the
    same discipline in both settings.
\end{itemize}
